\section{Reinforcement Learning Methodology and Experimental Protocol}
\label{sec:methods}

\subsection{Data Preparation and Leak-Free Temporal Alignment}
\label{subsec:data_prep}

To ensure strict experimental continuity and enable direct performance comparisons between static collaborative filtering and sequential decision-making, we preserve the exact data preprocessing pipeline established in Workshop~1 on the MovieLens~32M benchmark~\cite{harper2015movielens}. Using a fixed random seed ($\text{seed} = 42$) for full reproducibility, $15{,}000$ complete user histories are sampled without replacement from the raw interaction logs, followed by iterative $10$-core filtering that retains only users and movies with at least ten recorded interactions. This procedure yields a dense, high-signal corpus of $2{,}368{,}980$ explicit ratings (on a $0.5$ to $5.0$ half-star scale) spanning $14{,}998$ active users and $11{,}589$ distinct movies. Rather than partitioning ratings via random cross-validation---which violates causality by allowing future interactions to inform past predictions---all interactions are sorted globally by Unix timestamp and partitioned chronologically into an $80\%$ training set ($\mathcal{D}_{\text{train}}$, covering interactions from $1996$ through $2018$), a $10\%$ validation set ($\mathcal{D}_{\text{val}}$), and a $10\%$ test set ($\mathcal{D}_{\text{test}}$).

A critical methodological requirement in offline reinforcement learning for recommender systems is the prevention of temporal data leakage between representation learning and sequential policy evaluation. In our pipeline, the foundational latent representations---specifically the global rating mean $\mu \in \mathbb{R}$, user biases $b_u \in \mathbb{R}$, item biases $b_i \in \mathbb{R}$, and $20$-dimensional latent factor vectors $\mathbf{p}_u, \mathbf{q}_i \in \mathbb{R}^{20}$---are estimated via regularized FunkSVD matrix factorization~\cite{koren2009matrix} strictly on the historical training window $\mathcal{D}_{\text{train}}$ ($1996$--$2018$). Formally, these base parameters are obtained by minimizing the regularized squared error exclusively over training observations:
\begin{equation}
\label{eq:funksvd_train}
\resizebox{0.89\columnwidth}{!}{$\displaystyle \min_{\mathbf{p}_*, \mathbf{q}_*, b_*} \sum_{(u, i, r_{ui}) \in \mathcal{D}_{\text{train}}} \left( r_{ui} - \hat{r}_{ui} \right)^2 + \lambda_{\text{SVD}} \left( b_u^2 + b_i^2 + \|\mathbf{p}_u\|_2^2 + \|\mathbf{q}_i\|_2^2 \right),$}
\end{equation}
where the static point-estimate rating prediction is defined as
\begin{equation}
\label{eq:funksvd_pred}
\hat{r}_{ui} = \mu + b_u + b_i + \mathbf{p}_u \cdot \mathbf{q}_i.
\end{equation}
Similarly, all catalog popularity counts $|\mathcal{R}_i^{\text{train}}|$, genre frequency profiles, and empirical rating distributions used by the environment simulator are frozen at the boundary of $\mathcal{D}_{\text{train}}$. The chronological validation ($\mathcal{D}_{\text{val}}$) and test ($\mathcal{D}_{\text{test}}$) partitions are held out entirely from factor estimation and used exclusively to instantiate unseen evaluation episodes and verify ground-truth user satisfaction in post-$2018$ sessions. This strict temporal boundary ensures that the reinforcement learning agent cannot exploit future popularity surges or post-split preference drift when constructing state transitions or scoring candidate movies.

\subsection{Dueling Double Deep Q-Network (D3QN) Architecture}
\label{subsec:d3qn_arch}

Within the formulated Markov Decision Process $(\mathcal{S}, \mathcal{A}, \mathcal{P}, \mathcal{R}, \gamma)$~\cite{sutton2018reinforcement, sierra2026rl}, a stochastic or deterministic policy $\pi(a \mid s)$ maps states to action distributions to maximize the discounted return $G_t = \sum_{k=0}^{\infty} \gamma^k r_{t+k+1}$, evaluated via the state-value function $V^\pi(s) = \mathbb{E}[G_t \mid s_t = s]$ and action-value function $Q^\pi(s, a) = \mathbb{E}[G_t \mid s_t = s, a_t = a]$. Under classical model-free tabular Q-learning~\cite{watkins1992q, sierra2026rl}, the optimal action-value function $Q^*(s, a)$ satisfies the Bellman optimality equation:
\begin{equation}
\label{eq:bellman_optimality}
\resizebox{0.89\columnwidth}{!}{$\displaystyle Q^*(s_t, a_t) = \mathbb{E}_{s_{t+1} \sim \mathcal{P}(\cdot \mid s_t, a_t)} \left[ r_{t+1} + \gamma \max_{a' \in \mathcal{A}_{\text{valid}}(s_{t+1})} Q^*(s_{t+1}, a') \;\middle|\; s_t, a_t \right],$}
\end{equation}
which is approximated iteratively via the off-policy temporal-difference (TD) update rule $Q(s_t, a_t) \leftarrow Q(s_t, a_t) + \alpha\,\delta_t$ with TD error $\delta_t = r_{t+1} + \gamma \max_{a'} Q(s_{t+1}, a') - Q(s_t, a_t)$. Unlike on-policy SARSA ($Q(s_t, a_t) \leftarrow Q(s_t, a_t) + \alpha [r_{t+1} + \gamma Q(s_{t+1}, a_{t+1}) - Q(s_t, a_t)]$), which updates conservatively along the exploratory behavior policy~\cite{sutton2018reinforcement, sierra2026rl}, off-policy Q-learning directly estimates the greedy optimal policy regardless of the behavior policy and enables sample reuse via experience replay. However, tabular Q-learning fails in our sequential movie recommendation setting because the state space $\mathcal{S} \subset \mathbb{R}^{250}$ is continuous and high-dimensional, combining static user embeddings, dynamic session contexts, candidate item feature summaries, and session progress indicators; a lookup table of size $|\mathcal{S}| \times |\mathcal{A}|$ cannot store uncountably infinite floating-point states or generalize across unseen user sessions, necessitating neural function approximation $Q(s, a) \approx Q(s, a; \theta)$~\cite{mnih2015human, sierra2026rl}.

To achieve smooth generalization across $\mathcal{S} \subset \mathbb{R}^{250}$, we approximate the action-value function using a parameterized deep neural network $Q_\theta(s, a)$ built upon the Dueling Network Architecture~\cite{wang2016dueling}. Rather than mapping state representations directly to $|\mathcal{A}| = 50$ Q-values through a single monolithic feedforward path, the network first processes the input state vector $s \in \mathbb{R}^{250}$ through a shared two-layer multilayer perceptron (MLP) trunk equipped with Layer Normalization ($\text{LayerNorm}$) and rectified linear unit ($\text{ReLU}$) activations ($\mathbb{R}^{250} \to \mathbb{R}^{256} \to \mathbb{R}^{256}$):
\begin{align}
\label{eq:mlp_trunk_1}
\mathbf{h}^{(1)} &= \text{ReLU}\!\left(\text{LayerNorm}\!\left(\mathbf{W}^{(1)} s + \mathbf{b}^{(1)}\right)\right), \\
\label{eq:mlp_trunk_2}
\mathbf{h}^{(2)} &= \text{ReLU}\!\left(\text{LayerNorm}\!\left(\mathbf{W}^{(2)} \mathbf{h}^{(1)} + \mathbf{b}^{(2)}\right)\right),
\end{align}
where $\mathbf{W}^{(1)} \in \mathbb{R}^{256 \times 250}$, $\mathbf{W}^{(2)} \in \mathbb{R}^{256 \times 256}$, and $\mathbf{b}^{(1)}, \mathbf{b}^{(2)} \in \mathbb{R}^{256}$. Incorporating Layer Normalization within the shared trunk normalizes feature activations across the $256$ hidden units for each individual state sample, preventing covariate shift across heterogeneous state blocks (such as unnormalized latent dot products versus bounded genre frequencies) and stabilizing gradient flow during bootstrapping.

Following the shared representation $\mathbf{h}^{(2)} \in \mathbb{R}^{256}$, the architecture bifurcates into two specialized computational streams: a scalar state-value stream $V_\theta(s) \in \mathbb{R}$ and a $50$-dimensional action-advantage stream $A_\theta(s, a) \in \mathbb{R}^{50}$. Each stream is implemented as a dedicated two-layer head ($\mathbb{R}^{256} \to \mathbb{R}^{128} \to \mathbb{R}^{1}$ for $V_\theta$ and $\mathbb{R}^{256} \to \mathbb{R}^{128} \to \mathbb{R}^{50}$ for $A_\theta$) with intermediate ReLU activations. The two streams are subsequently recombined via the identifiability-preserving mean-subtraction aggregation layer~\cite{wang2016dueling}:
\begin{equation}
\label{eq:dueling_aggregation}
Q_\theta(s, a) = V_\theta(s) + \left( A_\theta(s, a) - \frac{1}{|\mathcal{A}|}\sum_{a'=0}^{|\mathcal{A}|-1} A_\theta(s, a') \right),
\end{equation}
where $|\mathcal{A}| = 50$ is the size of the candidate slate. Subtracting the mean advantage across all candidate slots forces the advantage stream to have zero mean at every state, resolving the unidentifiability issue wherein a constant scalar could otherwise be added to $V_\theta(s)$ and subtracted from $A_\theta(s, a)$ without altering $Q_\theta(s, a)$, while preserving the exact relative ordering of candidate actions $\arg\max_{a} Q_\theta(s, a)$.

This structural decomposition into $V_\theta(s)$ and $A_\theta(s, a)$ is particularly effective for sequential movie recommendation. In personalized recommendation sessions, the baseline expected return is heavily governed by the user's current engagement mood, historical rating generosity $b_u$, and remaining session horizon---properties captured globally by the state-value stream $V_\theta(s)$ independent of any single movie choice. Conversely, the advantage stream $A_\theta(s, a)$ isolates the relative marginal benefit of recommending specific candidate movie $a$ over the other $49$ candidates in the slate. In standard single-stream Q-networks, a single transition $(s_t, a_t, r_t, s_{t+1})$ updates only the output weight row corresponding to the executed action $a_t$, leaving the parameters of the remaining $49$ candidate slots untouched. Under the dueling formulation of~\eqref{eq:dueling_aggregation}, every temporal-difference backup directly updates $V_\theta(s_t)$, allowing the agent to rapidly learn accurate state valuations across all $50$ candidate actions simultaneously even when action execution is sparse.

\subsection{Double Q-Learning, Experience Replay, and Training Stability}
\label{subsec:ddqn_training}

While combining deep function approximation $Q(s, a; \theta)$ with Q-learning via the classical DQN squared TD loss $L(\theta) = \mathbb{E}[(r + \gamma \max_{a'} Q(s', a'; \theta^-) - Q(s, a; \theta))^2]$ enables scalable generalization~\cite{mnih2015human, sierra2026rl}, standard Deep Q-Networks (DQN) suffer from systematic maximization bias~\cite{van2016deep}. Because standard DQN uses the same target network $\theta^-$ both to select and evaluate the maximizing next action, positive estimation noise across the $50$ candidate slots is amplified by the $\max$ operator, resulting in overly optimistic Q-value estimates. To eliminate this upward bias, we employ Double Q-Learning (DDQN)~\cite{van2016deep, mnih2015human}, which decouples action selection from action evaluation by maintaining an online network with weights $\theta$ alongside a periodic target network with weights $\theta^-$. Given a transition tuple $(s_t, a_t, r_t, s_{t+1}, d_t)$ with terminal indicator $d_t \in \{0, 1\}$, the DDQN target $y_t$ is computed by using the online network $\theta$ to identify the greedy action over the valid unvisited candidate set $\mathcal{A}_{\text{valid}}(s_{t+1})$ and evaluating that chosen action under the target network $\theta^-$:
\begin{equation}
\label{eq:ddqn_target}
\resizebox{0.89\columnwidth}{!}{$\displaystyle y_t = r_t + (1 - d_t)\,\gamma\, Q_{\theta^-}\!\left(s_{t+1},\; \arg\max_{a' \in \mathcal{A}_{\text{valid}}(s_{t+1})} Q_\theta(s_{t+1}, a')\right).$}
\end{equation}
Here, invalid-action masking restricts the maximization strictly to $\mathcal{A}_{\text{valid}}(s_{t+1}) \subseteq \{0, 1, \dots, 49\}$, the subset of candidate slots not yet recommended during the ongoing episode. By masking already-recommended items ($Q_\theta(s, a) \leftarrow -\infty$ for $a \notin \mathcal{A}_{\text{valid}}(s)$) during both target computation and policy execution, the agent is guaranteed never to recommend duplicate movies within a session or bootstrap from invalid repetition states.

To train the online parameters $\theta$ robustly in the presence of noisy rating feedback and outlier temporal-difference errors, we minimize the Smooth $L_1$ (Huber) loss augmented with $L_2$ weight regularization over uniformly sampled mini-batches $B \sim \mathcal{D}$:
\begin{equation}
\label{eq:huber_objective}
\resizebox{0.89\columnwidth}{!}{$\displaystyle \mathcal{L}(\theta) = \frac{1}{|B|} \sum_{(s_t, a_t, r_t, s_{t+1}, d_t) \in B} \ell_{\text{Huber}}\!\left(Q_\theta(s_t, a_t) - y_t\right) + \lambda_{L_2} \|\theta\|_2^2,$}
\end{equation}
where $\lambda_{L_2} \ge 0$ is the $L_2$ regularization coefficient and the element-wise Huber penalty $\ell_{\text{Huber}}(\delta)$ for temporal-difference error $\delta = Q_\theta(s_t, a_t) - y_t$ transitions smoothly from quadratic to linear growth at $|\delta| = 1$:
\begin{equation}
\label{eq:huber_piecewise}
\ell_{\text{Huber}}(\delta) = 
\begin{cases} 
\frac{1}{2}\delta^2, & \text{if } |\delta| < 1, \\[4pt]
|\delta| - \frac{1}{2}, & \text{otherwise.}
\end{cases}
\end{equation}
Four complementary stabilization mechanisms prevent oscillatory divergence during optimization. First, we maintain a cyclic Experience Replay buffer $\mathcal{D}$ of capacity $N_{\text{replay}} = 50{,}000$ transitions and draw uniform random mini-batches of size $|B| = 128$ at each training step; randomizing transitions across thousands of historical user sessions breaks the strong temporal autocorrelation inherent to sequential steps within a single user episode~\cite{mnih2015human}. Second, the target network parameters $\theta^-$ are held fixed across training steps and synchronized periodically (or via Polyak averaging) with the online weights $\theta$, stabilizing the moving regression target $y_t$ in~\eqref{eq:ddqn_target}. Third, strict invalid-action masking ensures gradients are never corrupted by infeasible action backups. Fourth, prior to each Adam parameter update, we enforce global gradient norm clipping such that $\|\nabla_\theta \mathcal{L}(\theta)\|_2 \le 1.0$, bounding the magnitude of parameter steps when encountering unexpected high-reward transitions.

During training interaction, the agent balances candidate exploration and policy exploitation using an exponentially decaying $\epsilon$-greedy schedule restricted to valid unvisited candidates $\mathcal{A}_{\text{valid}}(s_t)$. At episode $t$, the exploration rate $\epsilon_t$ is governed by
\begin{equation}
\label{eq:epsilon_schedule}
\epsilon_t = \max\!\left(\epsilon_{\min},\; \epsilon_0 \cdot \delta_\epsilon^t\right),
\end{equation}
where $\epsilon_0 = 1.0$ denotes the initial pure-exploration rate, $\epsilon_{\min} = 0.05$ enforces a $5\%$ persistent exploration floor, and $\delta_\epsilon \in (0, 1)$ is the per-episode multiplicative decay factor. With probability $\epsilon_t$, the agent samples an action $a_t$ uniformly at random from $\mathcal{A}_{\text{valid}}(s_t)$; with probability $1 - \epsilon_t$, it selects the greedy action $a_t = \arg\max_{a \in \mathcal{A}_{\text{valid}}(s_t)} Q_\theta(s_t, a)$.

\subsection{Baseline Policies and Hyperparameter Search Protocol}
\label{subsec:baselines_tuning}

To rigorously quantify the value added by deep sequential planning and dueling value decomposition, we evaluate the trained D3QN agent against five deterministic and stochastic baseline policies executed on the exact same held-out test environment and candidate slates. Every baseline adheres to the identical invalid-action masking constraint, selecting actions exclusively from $\mathcal{A}_{\text{valid}}(s_t)$:

First, the \emph{Random Policy} ($\pi_{\text{rand}}$) selects an unvisited candidate uniformly at random at every step, $a_t \sim \text{Uniform}(\mathcal{A}_{\text{valid}}(s_t))$, establishing the lower-bound performance floor and characterizing the intrinsic reward density of the candidate generation pool. Second, the \emph{Popularity Heuristic} ($\pi_{\text{pop}}$) greedily selects the unvisited candidate with the highest historical interaction volume in the training corpus, $a_t = \arg\max_{a \in \mathcal{A}_{\text{valid}}(s_t)} |\mathcal{R}_{i(a)}^{\text{train}}|$. As established in Workshop~1, raw catalog popularity serves as a remarkably strong non-personalized ranking baseline due to heavy-tailed consumption patterns in MovieLens~32M, though it suffers from severe popularity bias and zero catalog novelty. Third, the \emph{Bias Baseline} ($\pi_{\text{bias}}$) ranks candidates using the non-personalized global and marginal rating offsets learned in Workshop~1, selecting $a_t = \arg\max_{a \in \mathcal{A}_{\text{valid}}(s_t)} (\mu + b_u + b_{i(a)})$, which captures baseline item quality without latent user-item interaction. Fourth, the \emph{Greedy FunkSVD Policy} ($\pi_{\text{SVD}}$, $\gamma = 0$, static) directly deploys the Workshop~1 matrix factorization model~\cite{koren2009matrix} as a sequential recommender by greedily selecting the candidate with the highest static predicted rating from~\eqref{eq:funksvd_pred}, $a_t = \arg\max_{a \in \mathcal{A}_{\text{valid}}(s_t)} (\mu + b_u + b_{i(a)} + \mathbf{p}_u \cdot \mathbf{q}_{i(a)})$. Because $\pi_{\text{SVD}}$ operates with a static user profile $\mathbf{p}_u$, zero discount horizon ($\gamma = 0$), and no awareness of novelty shaping or session fatigue, it directly measures the limitations of applying a static collaborative filter repeatedly across a session. Fifth, the \emph{Myopic Contextual Bandit} ($\pi_{\text{bandit}}$, $\gamma = 0$) incorporates dynamic intra-session preference updates and immediate novelty-shaped reward estimation at step $t$, but sets the reinforcement learning discount factor to $\gamma = 0$. By optimizing solely for immediate single-step reward without multi-step Bellman backups, $\pi_{\text{bandit}}$ serves as a direct ablation of the long-horizon sequential planning mechanism ($\gamma > 0$) central to our D3QN agent.

Finally, because deep reinforcement learning exhibits sensitivity to optimization and exploration hyperparameters~\cite{sutton2018reinforcement, mnih2015human}, we conduct a structured hyperparameter search across $8$ distinct D3QN agent configurations. The experimental grid systematically varies four core hyperparameters: the Adam optimizer learning rate $\eta \in \{10^{-3}, 5\times 10^{-4}, 10^{-4}\}$, the Bellman discount factor $\gamma \in \{0.0, 0.85, 0.95, 0.99\}$ controlling the effective planning horizon, the per-episode exploration decay factor $\delta_\epsilon \in \{0.985, 0.995, 0.998\}$ governing the transition speed from uniform exploration to greedy exploitation, and the $L_2$ weight decay coefficient $\lambda_{L_2} \in \{0, 10^{-5}, 10^{-4}\}$ regularizing the dueling network weights. All eight configurations are trained under identical random seeds, replay buffer capacities ($50{,}000$), and mini-batch sizes ($128$). Throughout training, agent checkpoints are evaluated periodically on validation episodes drawn from $\mathcal{D}_{\text{val}}$ using a purely greedy evaluation policy ($\epsilon = 0$). The configuration and training epoch achieving the highest exponentially smoothed validation return are saved to \texttt{checkpoints/agent\_best.pt} and locked for the final, single-shot comparative evaluation on the unseen test environment $\mathcal{D}_{\text{test}}$.
